% Section 7.1, from lectures/L07/appendix/a_avr_core.md.

\section{Mikrodatorn ATmega328P}
\label{c7:app:a}

\subsection{Från blockschema till en verklig krets}
\label{c7:a1}
I kapitel~\ref{c6:ch} byggdes mikrodatorn upp på blockschemanivå: en CPU med ALU, register och
styrenhet, ett programminne, ett dataminne, in- och utportar, och en klocka som driver alltihop.
Från och med nu arbetar vi med en verklig sådan krets: \textbf{ATmega328P}, som sitter på ett
Arduino Uno-kort.

ATmega328P är en \textbf{mikrokontroller}: en hel mikrodator på ett enda chip. Allt i
blockschemat från kapitel~\ref{c6:ch} finns inuti den:

\begin{booktable}{@{}p{8.6em}L@{}}
  \thead{Del} & \thead{I ATmega328P} \\
  \midrule
  CPU & En 8-bitars processor av typen AVR. Den räknar med 8 bitar, en byte, i taget. \\
  Register & 32 st, \code{r0}--\code{r31}, med 8 bitar vardera. \\
  Programminne & 32~kB flashminne. Programmet ligger kvar när strömmen bryts. \\
  Dataminne & 2~kB SRAM, 2048 byte. Innehållet försvinner när strömmen bryts. \\
  In- och utportar & Port B, C och D, som styr kretsens stift. \\
  Klocka & 16~MHz på ett Arduino Uno-kort: 16 miljoner klockcykler i sekunden. \\
\end{booktable}

Det finns inget operativsystem, inga fönster och ingen skärm. När kretsen får ström börjar den
köra det program som ligger i programminnet, instruktion för instruktion, och det fortsätter den
med så länge strömmen är på. Det är \textbf{ditt} program, och ingenting annat.

Två egenskaper hos maskinen är värda att känna till från början, eftersom de styr hur alla program
i boken ser ut:

\textbf{All räkning sker i registren.} Det finns ingen instruktion som adderar två byte i minnet
med varandra. Värdena hämtas till register, räknas där, och resultatet lämnas i ett register. Det
är därför bokens program handlar så mycket om \code{r16}, \code{r17} och deras grannar.

\textbf{Program och data ligger i olika minnen.} Programmet ligger i flashminnet och variablerna i
SRAM, och de två är skilda minnen med egna adresser (\secref{c7:a5}). Den uppdelningen kallas
\textbf{Harvardarkitektur}, och den känns igen från kapitel~\ref{c6:ch}.

\subsection{Registren}
\label{c7:a2}
Registren är processorns arbetsbänk: 32 st, vart och ett med plats för en byte, och med namnen
\code{r0}--\code{r31}. Allt processorn räknar på passerar genom dem.

\bookfigure{avr_register_file}{De 32 registren \code{r0}--\code{r31}.}{c7:fig:register_file}

De flesta instruktioner fungerar med vilket register som helst. Men en viktig grupp gör det inte:
\textbf{bara \code{r16}--\code{r31} kan laddas med en konstant.} Instruktionen \code{ldi r16, 25}
(\emph{load immediate}, ladda ett fast värde) är tillåten; \code{ldi r5, 25} är det inte, och
assemblern säger ifrån med ett felmeddelande. Samma begränsning gäller de andra instruktionerna som
tar en konstant, som \code{subi}, \code{andi}, \code{ori} och \code{cpi}, som du möter i
kapitel~\ref{c8:ch} och \ref{c9:ch}.

Förklaringen finns i hur instruktionen är kodad, och den står i \secref{c7:a6}. Den praktiska
regeln är enkel: \textbf{använd \code{r16}--\code{r25} som arbetsregister.} Det gör alla program i
boken, och då behöver du aldrig fundera på begränsningen. Registren \code{r26}--\code{r31} bildar
tre 16-bitars par, \code{X}, \code{Y} och \code{Z}, som används för att peka ut adresser i minnet;
dem låter vi vara i den här boken.

\subsection{Statusregistret SREG}
\label{c7:a3}
Utöver de 32 registren finns ett register som inte räknar själv, utan som \textbf{minns något om
den senaste uträkningen}: statusregistret \code{SREG}.

\bookfigure{avr_sreg}{Statusregistret \code{SREG} och dess åtta flaggor.}{c7:fig:sreg}

Varje bit i SREG kallas en \textbf{flagga}. Efter en addition kan processorn till exempel komma
ihåg om resultatet blev noll (flaggan \code{Z}, \emph{zero}) eller om det blev en minnessiffra ut
ur bit~7 (flaggan \code{C}, \emph{carry}). Flaggorna är det som gör att ett program kan fatta
beslut: \enquote{om resultatet blev noll, hoppa hit}. Hur varje flagga sätts är ämnet för
kapitel~\ref{c8:ch}, och hur hoppen läser dem är ämnet för kapitel~\ref{c9:ch}.

\subsection{Programräknaren}
\label{c7:a4}
Processorn måste hela tiden veta vilken instruktion som står på tur. Det håller ett särskilt
register reda på: \textbf{programräknaren}, \emph{program counter} eller \code{PC}. Den innehåller
adressen i programminnet till nästa instruktion.

Efter varje instruktion räknas \code{PC} upp till nästa instruktion, precis som räknaren i
\secref{c6:a7}. Det är därför ett program körs uppifrån och ned. Den enda gången programräknaren
gör något annat är när en instruktion \textbf{hoppar}: då skrivs en ny adress in i \code{PC}, och
programmet fortsätter där.

När kretsen startar, eller återställs med resetknappen, sätts \code{PC} till \code{0x0000}. Det är
därför första raden i varje program i boken ligger på adress 0 (se \secref{c7:b6}).

I simulatorn syns programräknaren som den gula pilen i editorn, och som värdet \emph{Program
Counter} i fönstret Processor Status.

\subsection{Tre minnen}
\label{c7:a5}
ATmega328P har tre minnen, och de har var sina adresser. Adress \code{0x0100} betyder alltså olika
saker beroende på vilket minne det gäller.

\bookfigure{avr_memory_spaces}{ATmega328P:s tre minnen, med sina adresser.}{c7:fig:memory_spaces}

\textbf{Programminnet} (flash, 32~kB) innehåller programmet: instruktionerna, en efter en. Det
behåller innehållet utan ström, vilket är varför ett Arduino-kort kör sitt program direkt när det
ansluts. Programminnet adresseras i \textbf{ord} om 16 bitar, eftersom nästan varje instruktion är
16 bitar lång.

\textbf{Dataminnet} (SRAM, 2~kB) innehåller det som ändras medan programmet kör. Det börjar på
adress \code{0x0100} och slutar på \code{0x08FF}, en adress som har ett eget namn, \code{RAMEND}.
Längst upp i dataminnet ligger \textbf{stacken}, som kapitel~\ref{c10:ch} handlar om. Under SRAM,
på de lägsta adresserna, ligger två andra saker som också räknas till dataminnet: de 32 registren
och \textbf{I/O-registren}.

\textbf{I/O-registren} är de register som styr kretsens delar utanför processorn. \code{DDRB},
\code{PORTB} och \code{PINB} styr port B, och \code{SREG} är också ett I/O-register. Ett
I/O-register läses och skrivs med egna instruktioner, \code{in} och \code{out}, som du möter i
kapitel~\ref{c8:ch}. Namnen, som \code{PORTB}, får du från filen \code{m328Pdef.inc} (se
\secref{c7:b3}).

\textbf{EEPROM} (1~kB) är ett litet minne som, liksom flash, tål att strömmen bryts, men som
programmet självt kan skriva i. Boken använder det inte.

\subsection{Vad en instruktion är}
\label{c7:a6}
En instruktion är, för processorn, bara ett tal: 16 bitar i programminnet. Det tal som står i
programminnet kallas \textbf{maskinkod}. Assemblerspråket, som du skriver, är ett sätt för
människor att skriva maskinkoden utan att räkna ut bitarna själva.

Ta instruktionen \code{ldi r16, 0x2A}. Den blir 16 bitar, och varje bit har en uppgift:

\bookfigure{avr_ldi_encoding}{Instruktionen \code{ldi r16, 0x2A} som sexton
bitar.}{c7:fig:ldi_encoding}

\begin{itemize}
  \item De fyra första bitarna, \code{1110}, säger att det är en \code{ldi}. De kallas
    \textbf{operationskoden}.
  \item De åtta bitarna märkta \code{K} är konstanten, \code{0x2A} = \code{0010 1010}, delad i två
    halvor.
  \item De fyra bitarna märkta \code{d} säger vilket register som ska laddas.
\end{itemize}

Här ser du också \textbf{varför \code{ldi} bara når \code{r16}--\code{r31}}. Fyra bitar kan bara
räkna från 0 till 15, alltså sexton olika register. Processorn lägger till 16 till talet i fältet,
så \code{0000} betyder \code{r16} och \code{1111} betyder \code{r31}. Det finns helt enkelt ingen
plats i instruktionen för att skriva \code{r5}.

Så \code{ldi r16, 0x2A} blir maskinkoden:

\begin{textdrawing}
1110  0010  0000  1010   =   0xE20A
 ldi  K7-4   r16  K3-0
\end{textdrawing}

Du kan kontrollera det själv i listfilen som Microchip Studio skriver när programmet byggs, se
\secref{c7:b7}. Nästan alla instruktioner i boken är ett ord, 16 bitar, långa.

\subsection{En instruktion per klockcykel}
\label{c7:a7}
Klockan på ett Arduino Uno-kort går med \textbf{16~MHz}. Det betyder 16 miljoner klockcykler per
sekund, och en klockcykel tar
\[
  T = \frac{1}{16\,000\,000\ \text{Hz}} = 62{,}5\ \text{ns}
\]

De flesta instruktioner tar \textbf{en klockcykel}. Några tar två, till exempel hopp; ett par tar
fler. Ett program med tio enkla instruktioner tar alltså bara 0,625 mikrosekunder. Det är så snabbt
att ingen människa kan se ett program köra, vilket är skälet till att vi stegar i simulatorn, och
till att programmen i kapitel~\ref{c9:ch} måste lägga in tidsfördröjningar för att en lysdiod ska
hinna synas.

Att antalet klockcykler är fast för varje instruktion är en ovanlig egenskap. Den gör att du kan
räkna ut exakt hur lång tid ett program tar, med papper och penna, innan du har kört det. Det gör
vi i kapitel~\ref{c9:ch}.

\subsection{Vad maskinen inte har}
\label{c7:a8}
En mikrokontroller är enkel, och det är värt att veta vad som saknas innan man letar efter det.

\textbf{Inget operativsystem att återvända till.} Ett program på en dator avslutas och lämnar
tillbaka kontrollen till operativsystemet. Här finns ingenting att lämna tillbaka till. Ett program
som \enquote{tar slut} fortsätter att köra det som råkar ligga efter det i programminnet. Därför
slutar varje program i boken med en loop som hoppar till sig själv, se \secref{c7:b6}.

\textbf{Inget skydd.} Ingenting hindrar ett program från att skriva i fel register eller hoppa till
fel ställe. Det blir inget felmeddelande när programmet körs, bara fel resultat. Därför är det så
viktigt att förutsäga vad varje instruktion ska göra, och kontrollera det i simulatorn.

\textbf{Ingen division.} Processorn kan addera, subtrahera och multiplicera, men inte dividera.
Division med 2, 4, 8 och så vidare görs genom att skifta bitarna (kapitel~\ref{c8:ch}).
